% Generated draft: requires mathematical and source review.
\documentclass[11pt]{article}
\usepackage[margin=25mm]{geometry}
\usepackage{fontspec}
\setmainfont{Latin Modern Roman}
\usepackage{amsmath,amssymb,mathtools,unicode-math}
\setmathfont{Latin Modern Math}
\usepackage{enumitem,needspace,xurl,hyperref}
\setlength{\parindent}{0pt}
\setlength{\parskip}{5pt}
\setlength{\emergencystretch}{3em}
\newcommand{\N}{\mathbb{N}}
\newcommand{\Z}{\mathbb{Z}}
\newcommand{\Q}{\mathbb{Q}}
\newcommand{\R}{\mathbb{R}}
\newcommand{\C}{\mathbb{C}}
\newcommand{\F}{\mathbb{F}}
\newcommand{\A}{\mathbb{A}}
\newcommand{\PP}{\mathbb{P}}
\newcommand{\E}{\mathbb{E}}
\newcommand{\eps}{\varepsilon}
\newcommand{\dd}{\,\mathrm{d}}
\begin{document}
\section*{UnsolvedMath ID: 2306116}
\textbf{Source problem code:} catalogue entry
\subsection{Definitions and mathematical statement}
Let \(D\subseteq\C\) be a domain containing \(0\). The source writes \(\Omega_t\) as the component of \(D\cap\{|z|\le t\}\) containing \(0\), and defines
\[
\omega_D(t)=\omega\bigl(0,\partial\Omega_t\cap\{|z|=t\},\Omega_t\bigr),
\]
where \(\omega(x,E,U)\) is harmonic measure of a boundary set \(E\) as seen from \(x\) in a domain \(U\). For instance, in the usual probabilistic interpretation it is the probability that planar Brownian motion exits \(U\) through \(E\).

The closed-disk intersection in the transcription need not be open. A natural domain-valued interpretation for \(t>0\) is instead
\[
U_t=\text{the component containing }0\text{ of }D\cap\{|z|<t\},
\qquad
\omega_D(t)=\omega\bigl(0,\partial U_t\cap\{|z|=t\},U_t\bigr).
\]
This correction must be checked; it is not silently identified with the literal set. The source declares a profile on \([0,\infty)\) but gives no harmonic-measure definition at \(t=0\).

Stephenson asks the following related questions about the map \(D\mapsto\omega_D\):
\begin{enumerate}
\item Characterize the functions \(w:[0,\infty)\to[0,1]\) realizable as \(\omega_D\).
\item If \(\omega_{D_1}\equiv\omega_{D_2}\), must the domains be the same up to rotation or reflection, or up to sets of capacity zero? What conclusions remain if equality is assumed only on a subinterval of \([0,1]\), as printed?
\item Given the full profile \(\omega_D\), reconstruct \(D\) in an appropriate equivalence sense.
\item Recover properties such as the connectivity of \(D\) from the profile; the source suggests beginning with circularly symmetric domains.
\end{enumerate}
Since a domain is already connected, ``connectivity'' here may concern simple/multiple connectivity or boundary-component structure, but the excerpt does not specify which. The capacity convention and the precise alternative equivalences in the uniqueness question are also unspecified. ``Circularly symmetric'' may mean rotational invariance or a circular-symmetrization class, which give different starting families. These are open formulation choices, not extra hypotheses imposed here. The supplied OPEN-TRIAGE summary reports open status in the 2018 edition and no definitive resolution located.
\subsection{Short English statement}
Characterize harmonic-measure radial profiles, determine how much they identify a planar domain, and reconstruct geometric or connectivity information. The truncated-domain and zero-radius definitions and the desired domain equivalence need clarification.
\subsection{Sources}
\begin{itemize}
\item Local UnsolvedMath registry, permanent ID \(2306116\), exact source code \texttt{catalogue entry}, ``Research Problems in Function Theory --- Problem 6.116.''
\item Hayman, \emph{Research Problems in Function Theory} (2018), Problem 6.116, attributed to K. Stephenson; \url{https://arxiv.org/abs/1809.07200}. Source-TeX extraction; recorded access 2026-07-29.
\item Supplied OPEN-TRIAGE summary, without a separate authoritative current-status citation.
\end{itemize}
\end{document}
